% Manuscript-ready fragment: raw harmonic powers at nonpositive spectral integers.
% Standard amsmath, amssymb, amsthm. Labels/citations have hp: prefix.
\section{Raw harmonic powers: polynomial finite parts and a centered cancellation}
\label{hp:sec:main}

The elementary-symmetric harmonic numerators of the height-one family do
not equal raw powers. For example, its second numerator is
$H_n^2-H_n^{(2)}$, whereas the family studied here contains $H_n^2$ itself.
Let
\begin{equation}\label{hp:def:F}
 H_n=\sum_{j=1}^n\frac1j,\qquad H_0=0,\qquad
 F_p(s,a)=\sum_{n=0}^{\infty}\frac{H_n^p}{(n+a)^s},
 \qquad p\ge1,\quad a>0,\quad \Re s>1.
\end{equation}
The inner harmonic numbers are unshifted; only the outer denominator
moves. All finite parts below mean the constant Laurent coefficient of
this meromorphic continuation in the variable $s$. They are not an
assignment of an ordinary value to the divergent defining series.

Dirichlet series of raw harmonic powers, their poles, and reductions to
Euler sums have already been studied, notably by Tin
\cite{hp:Tin2025}. The polylogarithm and shuffle ingredients are also
classical \cite{hp:IKZ2006}. The contribution developed here is the
continuous outer-shift normalization, its polynomial finite-part law at
every nonpositive integer, the exact cancellation at the centered shift,
and explicit low-weight identities with normalized parameter primitives.
No claim of priority is made for raw-power meromorphic continuation or
for the classical Euler sums used in the reductions.

\subsection{Every principal part has a finite Bernoulli formula}

Write $B_j(a)$ for the Bernoulli polynomial, with $B_1(a)=a-1/2$, and
introduce the formal series
\begin{equation}\label{hp:def:U}
 U(z,a)=-\sum_{j\ge1}\frac{B_j(a)}jz^j.
\end{equation}
Only finitely many coefficients are used; convergence of this formal
asymptotic series is neither required nor asserted.

\begin{theorem}[All poles and all principal coefficients]\label{hp:thm:poles}
The function $F_p(s,a)$ continues meromorphically to the entire $s$-plane,
locally uniformly in $a>0$. Its only possible poles are
$1,0,-1,-2,\ldots$. At $s=1$ its principal part is
\begin{equation}\label{hp:eq:principal-one}
 \sum_{r=0}^p\binom pr\frac{r!\gamma^{p-r}}{(s-1)^{r+1}}.
\end{equation}
At $s=-m$, $m\ge0$, put
\[
 F_p(-m+\varepsilon,a)
 =\sum_{h=1}^p\frac{R_{p,m,h}(a)}{\varepsilon^h}
       +C_{p,m}(a)+O(\varepsilon).
\]
Then every principal coefficient is the finite polynomial
\begin{equation}\label{hp:eq:allresidues}
 \boxed{\quad
 R_{p,m,h}(a)=\frac{p!}{(p-h+1)!}
       [z^{m+1}]\bigl(\gamma+U(z,a)\bigr)^{p-h+1},
 \qquad 1\le h\le p.\quad}
\end{equation}
In particular,
\begin{equation}\label{hp:eq:topresidue}
 R_{p,m,p}(a)=-\frac{p!}{m+1}B_{m+1}(a).
\end{equation}
Every $R_{p,m,h}$ has degree at most $m+1$ in $a$.
\end{theorem}

\begin{proof}
Put $x=n+a$. The classical shifted digamma expansion, obtained for
example from Euler--Maclaurin \cite{hp:DLMF}, is
\begin{equation}\label{hp:eq:digamma-expansion}
 H_n=\log x+\gamma-
       \sum_{j=1}^{J}\frac{B_j(a)}{jx^j}+O(x^{-J-1}).
\end{equation}
It holds uniformly when $a$ ranges over a compact positive interval;
the finite differentiated versions have the same uniformity. Raising a
finite truncation to the fixed power $p$ gives
\[
 H_n^p=\sum_{j=0}^{J}\sum_{r=0}^p
  \binom pr[z^j](\gamma+U(z,a))^{p-r}
                    x^{-j}(\log x)^r
       +O\bigl(x^{-J-1}(1+\log^p x)\bigr).
\]
Subtract these finitely many terms in the defining Dirichlet series.
The remainder is normally convergent in $\Re s>-J$ on compact
subsets. The subtracted sums are the corresponding linear combination
of $(-1)^r\partial_s^r\zeta(s+j,a)$. Arbitrarily large $J$ therefore
provides the meromorphic continuation and justifies every fixed
parameter derivative.

The only principal term of
$(-1)^r\partial_s^r\zeta(s+j,a)$ at $s=1-j$ is
$r!/(s+j-1)^{r+1}$. At $j=0$ this proves
\eqref{hp:eq:principal-one}. At $j=m+1$, set $r=h-1$ to obtain
\eqref{hp:eq:allresidues}. The case $r=p$ has no positive-degree
coefficient in $z$, so the pole order at a nonpositive integer is at
most $p$. Formula \eqref{hp:eq:topresidue} and the degree assertion
follow directly from the Bernoulli-polynomial coefficients.
\end{proof}

\begin{corollary}[Cancellation at every nonpositive even integer]
\label{hp:cor:centered}
For every $p\ge1$, the centered function $F_p(s,1/2)$ is holomorphic
at $s=0,-2,-4,\ldots$. Its only poles are at
$s=1,-1,-3,-5,\ldots$; their orders are exactly $p+1$ at $1$ and
exactly $p$ at each negative odd integer.
\end{corollary}
\begin{proof}
All odd Bernoulli polynomials vanish at $1/2$, including $B_1(1/2)$.
Thus $U(z,1/2)$ contains only even powers of $z$ and every coefficient
in \eqref{hp:eq:allresidues} vanishes when $m$ is even.
At $m=2r-1$, the top coefficient is
$-p!B_{2r}(1/2)/(2r)$, which is nonzero because
$B_{2r}(1/2)=(2^{1-2r}-1)B_{2r}\ne0$. The leading coefficient at
$s=1$ is $p!$ by \eqref{hp:eq:principal-one}.
\end{proof}

The cancellation concerns the entire Laurent principal part, not only
the residue. In the centered variable $x=n+1/2$ it is the exact
consequence of
$H_n=\log x+\gamma+1/(24x^2)-7/(960x^4)+\cdots$ having only even
inverse powers after its logarithm.

\subsection{A finite polylogarithm calculation determines every finite part}

For a positive-index word $\boldsymbol{k}=(k_1,\ldots,k_d)$, use
\[
 \operatorname{Li}_{\boldsymbol{k}}(z)
 =\sum_{n_1>\cdots>n_d\ge1}
       \frac{z^{n_1}}{n_1^{k_1}\cdots n_d^{k_d}}.
\]
Define the finite combination
\begin{equation}\label{hp:def:A}
 A_p(t)=\sum_{\substack{\boldsymbol{k}\text{ a composition of }p}}
     \frac{p!}{k_1!\cdots k_d!}
                  \operatorname{Li}_{\boldsymbol{k}}(e^{-t}).
\end{equation}
Ordering the distinct indices in the product $H_n^p$ proves the exact
identity
\begin{equation}\label{hp:eq:kernel}
 (1-e^{-t})\sum_{n\ge0}H_n^pe^{-nt}=A_p(t),\qquad
 \Gamma(s)F_p(s,a)=\int_0^\infty t^{s-1}K_{p,a}(t)\,dt,
 \quad K_{p,a}(t)=\frac{e^{-at}A_p(t)}{1-e^{-t}}.
\end{equation}
The integral is initially ordinary for $\Re s>1$. The first identity
is a finite stuffle decomposition: an ordered equality block of size
$k_j$ contributes $n_j^{-k_j}$ and its multiplicity is the displayed
multinomial coefficient.

For completeness, the local calculation has an explicit finite
recursion. With $L_{\boldsymbol{k}}(t)=
\operatorname{Li}_{\boldsymbol{k}}(e^{-t})$, set $L_{\varnothing}=1$.
Then
\begin{equation}\label{hp:eq:word-derivatives}
 \frac{d}{dt}L_{k_1,\ldots,k_d}(t)=
 \begin{cases}
  -L_{k_1-1,k_2,\ldots,k_d}(t),& k_1>1,\\
  -L_{k_2,\ldots,k_d}(t)/(e^t-1),& k_1=1.
 \end{cases}
\end{equation}
For a convergent word $k_1\ge2$, the constant is
$L_{\boldsymbol{k}}(0)=\zeta(\boldsymbol{k})$.
For a leading-one word, the integration constant is fixed by the
shuffle relations and $L_1(t)=-\log(1-e^{-t})$, whose constant
coefficient in powers of $\log t$ is zero. This is the usual shuffle
regularization. It is a finite reduction: multiplying a word with one
fewer initial $1$ by $\operatorname{Li}_1$ and shuffling its iterated
integral word isolates the desired initial-one word with a nonzero
integer multiplicity. All other terms have fewer initial $1$'s.
Induction reduces its constant to convergent multiple zeta values.

At a specified Taylor order, \eqref{hp:eq:word-derivatives} uses only
rational operations and finite integrations of logarithmic monomials.
Explicitly, for a polynomial $P$,
\begin{align}
 \int t^{-1}P(\log t)\,dt
   &=\int P(\ell)\,d\ell\big|_{\ell=\log t},\notag\\
 \int t^jP(\log t)\,dt
   &=t^{j+1}\sum_{q=0}^{\deg P}
           \frac{(-1)^qP^{(q)}(\log t)}{(j+1)^{q+1}},
                       \qquad j\ge0.\label{hp:eq:finite-integration}
\end{align}
This supplies an implementable calculation with fixed endpoint
constants, rather than an unspecified asymptotic coefficient.

\begin{theorem}[Polynomial finite parts at all nonpositive integers]
\label{hp:thm:finiteparts}
For fixed $p\ge1$ and $m\ge0$, compute the finite local expansion
\begin{equation}\label{hp:eq:local-kernel}
 K_{p,a}(t)=\sum_{j=-1}^{m}t^j
         \sum_{r=0}^p c_{p,j,r}(a)(\log t)^r
       +O\bigl(t^{m+1}(1+|\log t|^p)\bigr).
\end{equation}
The coefficients are provided by
\eqref{hp:def:A}--\eqref{hp:eq:finite-integration}; they are polynomials
in $a$ of degree at most $j+1$ with coefficients in the algebra of
convergent multiple zeta values. Define
\begin{equation}\label{hp:def:g}
 \frac1{\Gamma(-m+\varepsilon)}
      =\sum_{\ell\ge1}g_{m,\ell}\varepsilon^\ell.
\end{equation}
Then the exact finite part is
\begin{equation}\label{hp:eq:FP-all}
 \boxed{\quad C_{p,m}(a)=
   \sum_{r=0}^p(-1)^rr!\,c_{p,m,r}(a)g_{m,r+1}.\quad}
\end{equation}
In particular $C_{p,m}$ is a polynomial in $a$ of degree at most
$m+1$. Its coefficients lie in the algebra generated by Euler's
constant and convergent multiple zeta values. At the centered regular
points one has the stronger exact identity
\begin{equation}\label{hp:eq:center-value}
 F_p(-2r,1/2)=(2r)!\,c_{p,2r,0}(1/2).
\end{equation}
These centered values lie in the multiple-zeta algebra without requiring
Euler's constant.
\end{theorem}

\begin{proof}
Equations \eqref{hp:eq:word-derivatives} and
\eqref{hp:eq:finite-integration}, with the fixed shuffle constants,
construct the expansion by induction on the weight. The ordinary
convergent expansion of $1/(e^t-1)$ has a simple pole at zero and
radius $2\pi$. Finite integration of its analytic remainders preserves
a remainder bounded by a power of $t$ times a fixed logarithmic
polynomial. Consequently the displayed finite expansion is an actual
local expansion with the stated remainder, locally uniformly in $a$;
there are no negative powers below $t^{-1}$. Multiplication by
$e^{-at}$ proves the degree bound.

Subtract \eqref{hp:eq:local-kernel} inside the integral over $(0,1)$.
The remainder integral, together with the integral over $(1,\infty)$,
is holomorphic in $\Re s>-m-1$. Its fixed spectral derivatives are
obtained by inserting powers of $\log t$, by ordinary domination.
Thus
\begin{equation}\label{hp:eq:Mellin-continuation}
 \Gamma(s)F_p(s,a)=
 \sum_{j=-1}^{m}\sum_{r=0}^p
       \frac{(-1)^rr!c_{p,j,r}(a)}{(s+j)^{r+1}}+J_{p,m}(s,a),
\end{equation}
where $J_{p,m}$ is holomorphic near $s=-m$.
The zero of $1/\Gamma(s)$ at $s=-m$ kills the holomorphic remainder
and all terms with $j\ne m$ when taking the constant Laurent
coefficient. Only the finitely many $j=m$ terms remain, and their
product with \eqref{hp:def:g} is precisely \eqref{hp:eq:FP-all}.

For reference the reciprocal-Gamma coefficients are themselves finite
Bell expressions, obtained from
\begin{align}\label{hp:eq:gamma-generator}
 \frac1{\Gamma(-m+\varepsilon)}
 ={}&(-1)^m m!\,\varepsilon
 \exp\left((\gamma-H_m)\varepsilon\right.\notag\\
 &\left.\hspace{8mm}+
 \sum_{k\ge2}\frac{(-1)^{k+1}\zeta(k)-H_m^{(k)}}k
                                      \varepsilon^k\right).
\end{align}
The recurrence for Gamma proves this formula directly from the Taylor
series of $1/\Gamma(1+\varepsilon)$.
At a centered regular point, the highest possible logarithmic
coefficient in \eqref{hp:eq:Mellin-continuation} would produce a
nonzero pole. Descending triangularly, Corollary~\ref{hp:cor:centered}
therefore forces $c_{p,2r,j}(1/2)=0$ for every $j\ge1$. Since
$g_{2r,1}=(2r)!$, this proves \eqref{hp:eq:center-value}.
\end{proof}

\paragraph{Where the global information returns.}
The theorem evaluates the principal part and constant term. It does
not evaluate all spectral derivatives using the same finite data.
Already the coefficient of $\varepsilon$ receives the generally nonzero
contribution $(-1)^m m!J_{p,m}(-m,a)$ from the holomorphic remainder
in \eqref{hp:eq:Mellin-continuation}, together with the other regular
terms. In particular, differentiating a formula valid only for a
nonpositive integer does not determine the spectral derivative there.
This distinguishes the present finite identities from the higher
harmonic Stieltjes coefficients.

\subsection{The complete expansion through the finite part at zero}

Set
\begin{equation}\label{hp:def:P}
 P_p(L)=p![u^p]\exp(Lu)\Gamma(1+u)e^{\gamma u},\qquad P_0=1.
\end{equation}
Define ordinary convergent Euler sums
\begin{equation}\label{hp:def:E}
 E_{q,k}=\sum_{n\ge1}\frac{H_{n-1}^q}{n^k},
 \qquad q\ge0,\quad k\ge2.
\end{equation}
At $q=0$ the numerator is one. Every such sum is the finite
multiple-zeta combination
\begin{equation}\label{hp:eq:E-MZV}
 E_{q,k}=\sum_{\boldsymbol{r}\vDash q}
       \frac{q!}{r_1!\cdots r_d!}\zeta(k,r_1,\ldots,r_d),
\end{equation}
with the empty composition giving $E_{0,k}=\zeta(k)$.

\begin{theorem}[An exact all-degree zero-resonance identity]
\label{hp:thm:zero}
There are constants $Q_p$ with $Q_1=1/2$ and the finite recurrence
\begin{equation}\label{hp:eq:Q-recurrence}
 Q_p=-pQ_{p-1}+
 \frac12\sum_{k=2}^{p-1}(k-1)\binom p{k+1}E_{p-1-k,k},
                       \qquad p\ge2.
\end{equation}
For every $a>0$ one has
\begin{equation}\label{hp:eq:zero-expansion}
 \boxed{\quad F_p(s,a)=
   \left(\frac12-a\right)\frac{p!e^{\gamma s}}{s^p}
                 +Q_p+O(s).\quad}
\end{equation}
Consequently
\begin{equation}\label{hp:eq:zero-FP}
 C_{p,0}(a)=Q_p+\left(\frac12-a\right)\gamma^p,
                  \qquad F_p(0,1/2)=Q_p.
\end{equation}
Each $Q_p$ has a finite expression using convergent multiple zeta
values of weights at most $p-1$.
\end{theorem}

\begin{proof}
First the leading local polynomial of $A_p$ is $P_p(-\log t)$.
One way to fix it without a regularization guess is to use
\eqref{hp:eq:principal-one} in the Mellin identity. Multiplication by
$\Gamma(1+s)$ turns the elementary pole polynomial at $s=0$ into
exactly the coefficients of
$\Gamma(1+u)e^{\gamma u}$ in \eqref{hp:def:P}.
The word calculation proves existence of the local expansion, so this
identifies its leading polynomial uniquely.
At the centered shift
$e^{-t/2}/(1-e^{-t})=t^{-1}+O(t)$.
Regularity at $s=0$ then forces the coefficient of $t$ in $A_p$ to
have no logarithmic terms, by the same triangular argument as in the
previous theorem. Hence
\begin{equation}\label{hp:eq:A-first}
 A_p(t)=P_p(-\log t)+tQ_p+O(t^2(1+|\log t|^p)).
\end{equation}
The elementary formula for $A_1=-\log(1-e^{-t})$ gives $Q_1=1/2$.

To determine every other $Q_p$, put
\[
 M_{q,k}(t)=\sum_{n\ge1}\frac{H_{n-1}^q}{n^k}e^{-nt}.
\]
The exact increment $H_n^p-H_{n-1}^p$ gives
\[
 A_p=\sum_{k=1}^p\binom pkM_{p-k,k},\qquad
 A_p'=-\frac{pA_{p-1}}{e^t-1}
            -\sum_{k=2}^p\binom pkM_{p-k,k-1}.
\]
For $k\ge2$, $M_{q,k}(t)\to E_{q,k}$ by domination.
Using the preceding identity with $p-1$ also gives the complete
leading logarithmic polynomial
\[
 M_{p-2,1}(t)=\frac{1}{p-1}
 \left(P_{p-1}(-\log t)-
       \sum_{k=2}^{p-1}\binom{p-1}kE_{p-1-k,k}\right)+o(1).
\]
Take the coefficient of $t^0(\log t)^0$ in the derivative identity,
using \eqref{hp:eq:A-first} and
$(e^t-1)^{-1}=t^{-1}-1/2+O(t)$.
The two copies of $P_{p-1}$ cancel as full polynomials. This leaves
\[
 Q_p=-pQ_{p-1}+\sum_{k=2}^{p-1}
 \left(\frac p2\binom{p-1}k-\binom p{k+1}\right)E_{p-1-k,k}.
\]
The bracket is $(k-1)\binom p{k+1}/2$, proving the recurrence.
The local expansions already established justify differentiation and
coefficient comparison here.

It remains to show the particularly simple normalization in
\eqref{hp:eq:zero-expansion}. The coefficient of $t^0$ in the kernel is
$(1/2-a)P_p(-\log t)+Q_p$. If
$d(u)=\Gamma(1+u)e^{\gamma u}$, its $P_p$ contribution to the Mellin
continuation is
\[
 \frac{p!}{\Gamma(s)}[u^p]\frac{d(u)}{s-u}.
\]
Replace $d(u)$ by $d(s)$: the difference quotient is holomorphic in
$(s,u)$ near $(0,0)$ and its multiplication by $1/\Gamma(s)$ is
$O(s)$. The remaining term is
$p!e^{\gamma s}/s^p$. The constant $Q_p$ contributes
$Q_p/(s\Gamma(s))=Q_p+O(s)$. All other local and global terms are
$O(s)$ by the preceding theorem. This proves the result.
\end{proof}

The first five constants are
\begin{align}\label{hp:eq:Q-table}
 Q_1&=\frac12,& Q_2&=-1,\notag\\
 Q_3&=3+\frac12\zeta(2),&
 Q_4&=-12-2\zeta(2)+3\zeta(3),\notag\\
 Q_5&=60+10\zeta(2)-15\zeta(3)+\frac{33}{2}\zeta(4).
\end{align}
For example $E_{1,2}=\zeta(2,1)=\zeta(3)$ and
$E_{2,2}=2\zeta(2,1,1)+\zeta(2,2)=11\zeta(4)/4$.
These classical low-weight reductions follow from shuffle and stuffle
relations. Equation \eqref{hp:eq:zero-expansion} includes all negative
Laurent powers: for example the cubic case is
\[
 F_3(s,a)=\left(\frac12-a\right)
 \left(\frac6{s^3}+\frac{6\gamma}{s^2}+
                    \frac{3\gamma^2}{s}+\gamma^3\right)
         +3+\frac12\zeta(2)+O(s).
\]

\subsection{Explicit formulas at the next two spectral integers}

The finite word recursion gives the following local table. Here
$L=-\log t$, $z_j=\zeta(j)$, and the omitted remainder is
$O(t^4(1+|\log t|^p))$.
\begin{align}\label{hp:eq:A-table}
 A_1(t)={}&L+\frac t2-\frac{t^2}{24}+O(t^4),\notag\\
 A_2(t)={}&L^2+z_2-t-\frac{L}{12}t^2-\frac{t^3}{36}
                         +O(t^4(1+L^2)),\notag\\
 A_3(t)={}&L^3+3z_2L-2z_3+\left(3+\frac{z_2}{2}\right)t\notag\\
 &+\left(-\frac{L^2}{8}-\frac L4-\frac{z_2}{8}-\frac38\right)t^2
           +\frac5{72}t^3+O(t^4(1+|L|^3)),\notag\\
 A_4(t)={}&L^4+6z_2L^2-8z_3L+\frac{27}{2}z_4
             +(-12-2z_2+3z_3)t\notag\\
 &+\left(-\frac{L^3}{6}-\frac{L^2}{2}-\frac{z_2L}{2}-L
                 -\frac{z_2}{2}+\frac{z_3}{3}-\frac34\right)t^2
      +\left(-\frac{z_2}{18}-\frac7{27}\right)t^3
                     +O(t^4(1+|L|^4)).
\end{align}
This table is reproducible with \eqref{hp:eq:word-derivatives}.
For clarity the only weight-four endpoint reductions required are
\[
 \zeta(3,1)=\frac14\zeta(4),\qquad
 \zeta(2,2)=\frac34\zeta(4),\qquad
 \zeta(2,1,1)=\zeta(4).
\]
The leading-one regular constants follow by the indicated shuffle
elimination; for example they are $-5\zeta(4)/4$ for $(1,3)$,
$-3\zeta(4)$ for $(1,2,1)$, and $3\zeta(4)$ for $(1,1,2)$.
These constants, the two differentiation rules, and finite integration
fully determine every entry of \eqref{hp:eq:A-table}.

\begin{corollary}[Uniform square, cube, and fourth-power identities]
\label{hp:cor:minus-one}
For $1\le p\le4$, put $b=a-1/2$ and
\[
 T_1=0,\qquad T_2=-\frac1{12},\qquad
 T_3=\frac18,\qquad T_4=-\frac14.
\]
Then
\begin{equation}\label{hp:eq:minus-one-expansion}
 \boxed{\quad
 F_p(-1+\varepsilon,a)=p!e^{\gamma\varepsilon}
 \left(\frac{1/24-b^2/2}{\varepsilon^p}
               +\frac{b^2}{2\varepsilon^{p-1}}\right)
                  +bQ_p+T_p+O(\varepsilon).\quad}
\end{equation}
In particular,
\begin{equation}\label{hp:eq:minus-one-FP}
 C_{p,1}(a)=\frac{\gamma^p}{24}+T_p+bQ_p
                 +\frac{b^2}{2}(p\gamma^{p-1}-\gamma^p).
\end{equation}
The centered regular values at $s=-2$ are
\begin{align}\label{hp:eq:minus-two-table}
 F_1(-2,1/2)&=-\frac1{24},&
 F_2(-2,1/2)&=\frac1{36},\notag\\
 F_3(-2,1/2)&=-\frac19-\frac{\zeta(2)}{24},&
 F_4(-2,1/2)&=\frac{13}{27}+\frac{\zeta(2)}{18}
                                  -\frac{\zeta(3)}4.
\end{align}
\end{corollary}

\begin{proof}
Multiply \eqref{hp:eq:A-table} by
$e^{-at}/(1-e^{-t})$ and apply \eqref{hp:eq:FP-all} at $m=1$.
Together with the Bernoulli formula for the principal part this gives
\eqref{hp:eq:minus-one-expansion}. Equivalently, the centered calculation
uses $e^{-t/2}/(1-e^{-t})=t^{-1}-t/24+O(t^3)$ and the
reciprocal-Gamma expansion at $-1$; the exact parameter transport
proved next supplies its entire $b$-dependence.
For $m=2$, no logarithmic term survives at the centered shift, and
\eqref{hp:eq:center-value} says to multiply the constant coefficient
of $t^2$ by $2$. This coefficient is
$[t^3]A_p-Q_p/24$. Inserting the displayed table gives all four
values in \eqref{hp:eq:minus-two-table}.
\end{proof}

\subsection{Exact shift derivatives and normalized antiderivatives}

Let $R_{p,m}(a)=R_{p,m,1}(a)$ be the simple-pole coefficient in
\eqref{hp:eq:allresidues}. It is given without a polylogarithm
calculation by
\begin{equation}\label{hp:eq:residue-poly}
 R_{p,m}(a)=[z^{m+1}](\gamma+U(z,a))^p.
\end{equation}

\begin{theorem}[Parameter transport and finite primitives]
\label{hp:thm:transport}
For $m\ge1$,
\begin{equation}\label{hp:eq:transport}
 C_{p,m}'(a)=mC_{p,m-1}(a)-R_{p,m-1}(a).
\end{equation}
At $m=0$ the corresponding identity is $C_{p,0}'(a)=-\gamma^p$.
For every $m\ge0$ and positive $a,a_0$ one consequently has the
normalized antiderivative identity
\begin{equation}\label{hp:eq:primitive}
 \boxed{\quad
 \int_{a_0}^a C_{p,m}(x)\,dx=
 \frac{C_{p,m+1}(a)-C_{p,m+1}(a_0)
              +\int_{a_0}^aR_{p,m}(x)\,dx}{m+1}.\quad}
\end{equation}
The last integral is a finite polynomial integral explicitly specified
by \eqref{hp:eq:residue-poly}; it introduces no new special function
or undetermined normalization.
\end{theorem}
\begin{proof}
Normal convergence in the defining half-plane gives
$\partial_aF_p(s,a)=-sF_p(s+1,a)$, and the normally convergent
continuations extend this identity meromorphically. Take the constant
coefficient at $s=-m+\varepsilon$.
The factor $m-\varepsilon$ contributes
$mC_{p,m-1}-R_{p,m-1}$, proving \eqref{hp:eq:transport}.
At $m=0$, the residue at $s=1$ is $\gamma^p$ by
\eqref{hp:eq:principal-one}. Integrating the relation with $m+1$
gives \eqref{hp:eq:primitive} and fixes its additive constant.
\end{proof}

For example, substituting \eqref{hp:eq:zero-FP} and
$R_{p,0}(a)=p(1/2-a)\gamma^{p-1}$ into
\eqref{hp:eq:transport} proves the quadratic shift dependence
\eqref{hp:eq:minus-one-FP} from its centered value. This provides an
independent exact check on the $a$-polynomial obtained by Mellin
coefficient extraction.

\paragraph{Extension and remaining questions.}
A polynomial in finitely many generalized harmonic numbers
$H_n^{(r)}$ of positive integer orders $r$ admits the same finite
ordered-equality-block expansion.
Replacing \eqref{hp:def:A} by its finite multiple-polylogarithm
combination (including a constant when the polynomial has one)
therefore preserves the finite-part theorem, polynomial
shift dependence, and normalized parameter transport. Its pole orders
must be read from its own finite asymptotic expansion; the particular
centered cancellation for raw powers does not automatically extend to
arbitrary generalized-harmonic products.
Three concrete questions remain: determine how far the $Q_p$ reduce
within specified multiple-zeta coordinate spaces at higher weight;
produce compact all-degree formulas for the centered values
$F_p(-2r,1/2)$ beyond the word recursion; and identify combinations of
first regular spectral coefficients whose global Mellin remainders
cancel. The last question is distinct from the finite-part identities
proved here.
